% !TeX program = lualatex
% !TeX encoding = utf8
% !TeX spellcheck = uk_UA

\documentclass[]{beamer}
\usetheme{PlasmaPhysics}
\usepackage{PlasmaPhysics}
\usetikzlibrary{shadings}
\graphicspath{{pictures}}

%\NewDocumentCommand{\E}{}{\vec{E}}
%\NewDocumentCommand{\B}{}{\vec{B}}
%\NewDocumentCommand{\J}{}{\vec{j}}
%\NewDocumentCommand{\V}{}{\vec{v}}
%\NewDocumentCommand{\wk}{}{\vec{k}}
%
%\NewDocumentCommand{\divg}{m}{\vec{\nabla}\cdot{#1}}
%\NewDocumentCommand{\rot}{m}{\vec{\nabla}\times{#1}}
%\NewDocumentCommand{\grad}{m}{\vec{\nabla}#1}
%\NewDocumentCommand{\vecdot}{mm}{\left[#1\times#2\right]}
%\NewDocumentCommand{\scdot}{mm}{\left(#1\cdot#2\right)}
%\NewDocumentCommand{\dt}{m}{\frac{\partial #1}{\partial t}}
%\NewDocumentCommand{\ddt}{m}{\frac{\partial^2 #1}{\partial t^2}}

\usepackage{subcaption}
\usepackage{mathtools}
\setlength{\abovedisplayskip}{3pt}
\setlength{\belowdisplayskip}{3pt}
\usetikzlibrary{decorations.pathreplacing}
\newtcolorbox{mysidebyside}[1][]{
    notitle, sidebyside, blanker,
    sidebyside align=center seam,
    halign=center,
    lefthand width=.5\textwidth,
    #1
}



\graphicspath{{pictures/}}
\title[Лекції з фізики плазми]{\bfseries\huge Стійкість магнітогідродинамічних конфігурацій}
\subtitle{\bfseries Лекції з фізики плазми}
\author{Пономаренко С. М.}
\date{}
% https://www.youtube.com/watch?v=zv9Y7YmHyBk
\def\stext#1{\text{
        \makebox[0pt]{
            \tikz[baseline] {\node[text width=50pt, align=center, execute at begin node=\setlength{\baselineskip}{-0.5ex}, blue] {#1};}
        }
}}

\begin{document}

% --- Слайд 1 ---
\begin{frame}
	\titlepage
\end{frame}

\begin{frame}{МГД-нестійкості: фізичний зміст}{}
	\begin{block}{}\justifying
		Магнітогідродинаміка (МГД) описує плазму як провідну рідину в електромагнітному полі.
		Стаціонарні конфігурації такої системи можуть втрачати стійкість.
	\end{block}

	\begin{block}{Ідея нестійкості}
		Мала збурююча дія може:
		\begin{itemize}
			\item або затухати $\rightarrow$ система \textbf{стійка}
			\item або зростати $\rightarrow$ система \textbf{нестійка}
		\end{itemize}
	\end{block}

	\begin{block}{Фізичне джерело}
		\begin{itemize}
			\item градієнти тиску
			\item струми в плазмі
			\item конфігурація магнітного поля
		\end{itemize}
	\end{block}
\end{frame}

\begin{frame}{Лінійний аналіз стійкості}{}
	\begin{block}{Основна ідея}
		Розглядаємо мале збурення відносно рівноваги:
		\[
			\mathbf{X} = \mathbf{X}_0 + \tilde{\mathbf{X}},
		\]
		$\mathbf{X}_0$ --- стаціонарний стан, $\tilde{\mathbf{X}} \ll \mathbf{X}_0$ --- мале збурення.
	\end{block}

	\begin{block}{Лінеаризація}
		Підставляємо у рівняння МГД та залишаємо тільки перший порядок малості:
		\[
			\frac{\partial \tilde{\mathbf{X}}}{\partial t} = \hat{L} \tilde{\mathbf{X}}
		\]
	\end{block}

	\begin{block}{}
		Розв'язок шукаємо у вигляді:
		\(
		\tilde{\mathbf{X}} \sim e^{i(\mathbf{k}\cdot \mathbf{r} - \omega t)}
		\)
	\end{block}

	\begin{block}{Критерій стійкості}
		\begin{itemize}
			\item $\operatorname{Im}(\omega) > 0$ --- нестійкість (експоненційний ріст)
			\item $\operatorname{Im}(\omega) < 0$ --- затухання
		\end{itemize}
	\end{block}
\end{frame}

\begin{frame}{Нестійкість Релея--Тейлора: аналогія}{}
	\begin{block}{Класична гідродинаміка}
		\begin{itemize}
			\item Дві рідини з різною густиною:
			      \(
			      \rho_1 > \rho_2
			      \)
			\item Важка рідина (\(\rho_1\)) знаходиться \textbf{над} легкою (\(\rho_2\))
			\item Гравітація напрямлена вниз
		\end{itemize}
	\end{block}

	\begin{block}{Що відбувається?}
		\begin{itemize}
			\item Система "перевернута"
			\item Будь-яке збурення межі:
			      \begin{itemize}
				      \item важка рідина опускається
				      \item легка піднімається
			      \end{itemize}
			\item Збурення зростає $\rightarrow$ \textbf{нестійкість}
		\end{itemize}
	\end{block}

	\begin{block}{Аналогія в плазмі}
		\begin{itemize}
			\item Замість гравітації:
			      \[
				      \mathbf{g}_{\text{eff}} \sim -\frac{\nabla p}{\rho}
			      \]
			\item "Важке" = область з більшим тиском / густиною
			\item Якщо плазма "підперта неправильно" $\rightarrow$ виникає нестійкість
		\end{itemize}
	\end{block}
\end{frame}

\begin{frame}{Нестійкість Релея--Тейлора: ріст збурень}{}
	\begin{block}{Модель}
		\begin{itemize}
			\item Дві незмішувані рідини з густинами $\rho_1$, $\rho_2$
			\item Межа між ними збурюється:
			      \[
				      \eta \sim e^{i(kx - \omega t)}
			      \]
		\end{itemize}
	\end{block}

	\begin{block}{Результат}
		Швидкість росту збурення:
	\end{block}

	\begin{center}
		:contentReference[oaicite:0]{index=0}
	\end{center}

	\begin{block}{Позначення}
		\begin{itemize}
			\item $\gamma = \operatorname{Im}(\omega)$ --- інкремент росту
			\item $k$ --- хвильове число
			\item $A = \dfrac{\rho_1 - \rho_2}{\rho_1 + \rho_2}$ --- число Атвуда
		\end{itemize}
	\end{block}

	\begin{block}{Фізичний зміст}
		\begin{itemize}
			\item Якщо $A > 0$ (важке над легким) $\rightarrow$ нестійкість
			\item Чим більше $k$ $\rightarrow$ тим швидше росте збурення
		\end{itemize}
	\end{block}
\end{frame}

\end{document}